\documentclass[11pt,letterpaper]{article}
\usepackage[T1]{fontenc}
\usepackage[utf8]{inputenc}
\usepackage{lmodern,geometry,microtype,xcolor,enumitem,hyperref,fancyhdr}
\geometry{margin=0.78in}
\definecolor{ink}{HTML}{17324D}
\definecolor{accent}{HTML}{126A78}
\hypersetup{colorlinks=true,linkcolor=accent,urlcolor=accent,breaklinks=true}
\setlength{\parindent}{0pt}
\setlength{\emergencystretch}{3em}
\setlist[itemize]{leftmargin=*,itemsep=3pt,topsep=3pt}
\pagestyle{fancy}\fancyhf{}\fancyhead[L]{\small Architecture Documentation Suite}\fancyhead[R]{\small v1.0}\fancyfoot[R]{\thepage}\setlength{\headheight}{15pt}
\begin{document}
\thispagestyle{empty}
{\color{accent}\large TECHNICAL REFERENCE \hfill NOVEMBER 2026}\par\vspace{1.1cm}
{\color{ink}\LARGE\bfseries\raggedright Integrated Architecture Documentation Framework}\par\vspace{0.35cm}
{\Large V\&B + C4 + UML + PlantUML + ISO/IEC/IEEE 42010:2022 alignment}\par\vspace{0.55cm}
\rule{\textwidth}{0.7pt}\par\smallskip
{\small Architecture Documentation Suite | Release 1.0}\par\vspace{0.7cm}
\tableofcontents
\newpage
\section*{1. Framework roles and limits}\addcontentsline{toc}{section}{1. Framework roles and limits}
Views and Beyond determines what information to capture and how to package it. C4 supplies a zoomable hierarchy for static software-system diagrams and supporting dynamic/deployment diagrams. UML supplies defined modeling constructs for structural and behavioral descriptions. PlantUML is a text-based rendering notation/tool, not a complete architecture-description methodology. ISO/IEC/IEEE 42010:2022 defines architecture-description requirements and also addresses viewpoints, model kinds, description frameworks and languages. The full current normative text was not supplied; this framework is alignment-oriented and does not claim clause-level conformance.\par\smallskip
\section*{2. Common metamodel}\addcontentsline{toc}{section}{2. Common metamodel}
Use stable IDs: SYS for systems, E for elements, R for relations, C for concerns, ST for stakeholders, V for views, VP for viewpoints, MK for model kinds, ADR for decisions and EV for evidence. Use unique IDs within a baseline; preserve mappings across revisions. Every claim is tagged observed, inferred, proposed, or unknown. Observed claims require repository or runtime evidence; inferred claims state assumptions; proposed claims have an explicit design status. Documentation evidence cannot itself prove a runtime invariant.\par\smallskip
\section*{3. View selection and notation crosswalk}\addcontentsline{toc}{section}{3. View selection and notation crosswalk}
\begin{itemize}
\item Context viewpoint: system boundary, people and external systems; C4 system context; addresses scope, ownership and integration. Context may also appear inside a V\&B packet.
\item Runtime service viewpoint: applications, stores, interactions and protocols; C4 container diagram plus explicit C\&C connector semantics. A C4 container is an application or data store, not necessarily a Docker container.
\item Module viewpoint: implementation responsibilities and dependencies; UML packages, classes or dependency matrix; C4 component/code diagrams only when their chosen semantics match. Never assume C4 components and V\&B runtime components are equivalent.
\item Runtime component viewpoint: internal interacting units, ports and connectors; UML component or suitable C4 component diagram with a runtime contract. Distinguish static dependencies from calls.
\item Deployment viewpoint: allocation to execution nodes; C4 deployment or UML deployment; state environment, instance counts and network boundaries.
\item Behavior viewpoint: interactions and lifecycle; C4 dynamic for selected collaborations, UML sequence/activity/state for semantics requiring more detail.
\item Security viewpoint: assets, trust boundaries, identities, enforcement and data flows; DFD or typed PlantUML representation, supplemented by threat and control tables. This is a local extension, not a C4 level.
\item Data viewpoint: schema, lifecycle, cardinalities and ownership; UML class/entity model and data-flow mapping when required.
\end{itemize}
\section*{4. Viewpoint specification recipe}\addcontentsline{toc}{section}{4. Viewpoint specification recipe}
For each VP define purpose, concern set, intended stakeholders, model kinds, element/relationship grammar, allowed constraints, construction inputs, analysis procedures and consistency checks. For each model kind define allowed types, key properties, presentation conventions and verification procedure. Example VP-MOD uses package models whose arrow means implementation dependency; allowed constraints are no adapter imports from domain and no cyclic layer dependencies. Example VP-RUN uses directed requests with initiating component, protocol and timeout; it must not silently reuse the static arrow meaning.\par\smallskip
\section*{5. Correspondence rules}\addcontentsline{toc}{section}{5. Correspondence rules}
\begin{itemize}
\item CR-01: each internal runtime container corresponds to at least one implementation unit or has a documented external/generated origin.
\item CR-02: each deployed instance references a container/runtime element; each environment has its own mapping.
\item CR-03: each behavior message corresponds to a runtime connector or is explicitly marked as local/internal behavior.
\item CR-04: a trust-boundary crossing names an identity, authorization point, data classification and enforcement control or a gap.
\item CR-05: each high-priority concern has a selected view, evidence and review disposition; not-addressed concerns must remain visible.
\item CR-06: type mismatches, orphan IDs and stale baseline references become tracked inconsistencies, not automatically reconciled claims.
\end{itemize}
These are proposed project-level rules; tailor cardinalities to the system. Example: a single domain package may execute in API and worker processes, so the module-to-runtime mapping is many-to-many.\par\smallskip
\section*{6. Alignment register (non-normative)}\addcontentsline{toc}{section}{6. Alignment register (non-normative)}
Use a review matrix with fields: topic, applicable normative requirement reference (to be completed using licensed 2022 text), artifact location, evidence, reviewer and disposition. Candidate topics are AD identification and scope, stakeholders and concerns, viewpoint and model-kind conventions, architecture views/models, correspondence and inconsistencies, architecture decisions and rationale, and conformance assessment. This list is not represented as the standard's exhaustive requirements. The book's Epilogue E.1 relates V\&B to the earlier standard; do not copy its conformance checklist as a current audit.\par\smallskip
\section*{7. Diagram conventions}\addcontentsline{toc}{section}{7. Diagram conventions}
Use title, scope, baseline, key and stable element labels. Arrows must have verb semantics and direction, plus technology/protocol where meaningful. Use system boundaries and deployment boundaries differently. Avoid color-only semantics. Store .puml sources alongside the architecture description and render using a pinned, approved PlantUML version. The supplied diagrams use native PlantUML and C4-style labels rather than a network-loaded macro library. UML and C4 diagrams can coexist; do not call every C4 diagram UML.\par\smallskip
\section*{8. Documentation pipeline}\addcontentsline{toc}{section}{8. Documentation pipeline}
Select concerns and viewpoints; collect evidence; normalize catalogs; author presentations and supporting packets; map across views; conduct stakeholder review; publish baseline. Maintain the canonical model, not duplicated prose facts. On pull requests affecting architecture-sensitive paths, require either an AD update or a documented no-impact decision. Build diagrams offline; run structural checks and validate source citations. Human review remains essential for correctness and standards conformance.\par\smallskip
\section*{9. Source register}\addcontentsline{toc}{section}{9. Source register}
V\&B: Clements et al., chapters 1--11 and E.1. ISO public record: \url{https://www.iso.org/standard/74393.html}. C4 official hierarchy and notation: \url{https://c4model.com/diagrams} and \url{https://c4model.com/diagrams/notation}. PlantUML: \url{https://plantuml.com/}. See document 05 for the suite's local acceptance thresholds and document 07 for worked application.\par\smallskip
\section*{Bibliographic foundation}
Paul Clements et al., \emph{Documenting Software Architectures: Views and Beyond}, 2nd ed., Addison-Wesley/Pearson, copyright 2011, ISBN 978-0-321-55268-6. Section and chapter references refer to the supplied edition. Additional framework sources are identified in the applicable documents.\par
\end{document}
